\chapter{Chapter06 — Bent Functions}
%%% generated by the blueprint pipeline from FABL.Chapter06.Constructions.BentFunctions
\section{Overview}

Declarations of the Lean module \texttt{FABL.Chapter06.Constructions.BentFunctions}.

\section{Declarations}

\begin{definition}[Sign-valued function]
\lean{FABL.IsSignValued}
\label{FABL.IsSignValued}
\leanok
\uses{FABL.F₂Cube}
A function $f:(\mathbb F_2)^n\to\mathbb R$ is sign-valued if $|f(x)|=1$ for every
$x\in(\mathbb F_2)^n$.
\end{definition}

\begin{theorem}[Sign-valued real encoding]
\lean{FABL.isSignValued_realSignEncodedFunction}
\label{FABL.isSignValued_realSignEncodedFunction}
\leanok
\uses{FABL.F₂BooleanFunction, FABL.IsSignValued, FABL.realSignEncodedFunction, FABL.signEncode, FABL.signEncodedFunction, FABL.signValue_eq_neg_one_or_one}
\difficulty{4}
Let $n$ be a nonnegative integer and let $f:\mathbb F_2^n\to\mathbb F_2$. Encode each
value $f(x)$ as the real sign $(-1)^{f(x)}$; the resulting real-valued function has
absolute value $1$ at every $x\in\mathbb F_2^n$.
\end{theorem}

\begin{definition}[Bent function condition]
\lean{FABL.IsBent}
\statementsource{Analysis-of-Boolean-Functions-by-Ryan-ODonnell}{
  pdf-ffbe9c7b3306-p151-b005,
  pdf-ffbe9c7b3306-p170-b005
}
\label{FABL.IsBent}
\leanok
\uses{FABL.F₂Cube, FABL.vectorFourierCoeff}
A real-valued function $f$ on the $n$-dimensional Boolean cube is bent if every vector
Fourier coefficient of $f$ has absolute value $(2^{loor{n/2}})^{-1}$.
\end{definition}

\begin{theorem}[Parseval identity for vector Fourier coefficients]
\lean{FABL.sum_sq_vectorFourierCoeff_eq_one}
\label{FABL.sum_sq_vectorFourierCoeff_eq_one}
\leanok
\uses{FABL.F₂Cube, FABL.IsSignValued, FABL.vectorFourierCoeff, FABL.vector_plancherel}
\difficulty{4}
Let $n$ be a natural number and let $f$ be a real-valued function on the $n$-dimensional
vector space over $\\mathbb F_2$. Suppose $|f(x)|=1$ for every input $x$. Then the sum,
over all vectors $\\gamma$ in that space, of the squares of the Fourier coefficients
indexed by the support of $\\gamma$ is $1$.
\end{theorem}

\begin{theorem}[A large binary Fourier coefficient]
\lean{FABL.exists_inv_card_le_sq_vectorFourierCoeff}
\label{FABL.exists_inv_card_le_sq_vectorFourierCoeff}
\leanok
\uses{FABL.F₂Cube, FABL.IsSignValued, FABL.sum_sq_vectorFourierCoeff_eq_one, FABL.vectorFourierCoeff, FABL.𝔽₂}
\difficulty{5}
Let $n$ be a nonnegative integer, and let $f:(\\mathbb{F}_2)^n\\to\\mathbb{R}$ satisfy
$|f(x)|=1$ for every $x$. Then some binary Fourier coefficient of $f$ has square at
least $2^{-n}$.
\end{theorem}

\begin{theorem}[A large Fourier coefficient for sign-valued functions]
\lean{FABL.exists_bent_extremal_coefficient}
\label{FABL.exists_bent_extremal_coefficient}
\leanok
\uses{FABL.F₂Cube, FABL.IsSignValued, FABL.exists_inv_card_le_sq_vectorFourierCoeff, FABL.vectorFourierCoeff}
\difficulty{5}
Let $n$ be even, and let $f:\{0,1\}^n\to\mathbb R$ satisfy $|f(x)|=1$ for every $x$.
Then some Fourier coefficient of $f$ on the binary cube has absolute value at least
$2^{-n/2}$.
\end{theorem}

\begin{theorem}[Bentness of the mod-two inner-product function]
\lean{FABL.isBent_innerProductModTwo}
\proofsource{Analysis-of-Boolean-Functions-by-Ryan-ODonnell}{
  pdf-ffbe9c7b3306-p152-b006,
  pdf-ffbe9c7b3306-p152-b007
}
\label{FABL.isBent_innerProductModTwo}
\leanok
\uses{FABL.IsBent, FABL.abs_vectorFourierCoeff_innerProductModTwo, FABL.innerProductModTwo}
\difficulty{3}
For every natural number $n$, the function on pairs of $n$-bit vectors that maps $(x,y)$
to $(-1)^{\langle x,y\rangle}$, with the inner product computed modulo $2$, is bent on
the $2n$-dimensional binary cube: every Fourier coefficient has absolute value $2^{-n}$.
\end{theorem}

\begin{theorem}[The Inner-Product Sign Function]
\lean{FABL.isSignValued_innerProductModTwo}
\label{FABL.isSignValued_innerProductModTwo}
\leanok
\uses{FABL.IsSignValued, FABL.f₂CubeBlockEquiv, FABL.f₂DotProduct, FABL.innerProductModTwo, FABL.innerProductModTwo_joinF₂CubeBlocks, FABL.joinF₂CubeBlocks}
\difficulty{5}
For every natural number $n$, the function on pairs $x,y\\in\\mathbb F_2^n$ given by
$(-1)^{\\sum_{j=1}^n x_jy_j}$, with the exponent computed modulo $2$, has absolute value
$1$ at every pair.
\end{theorem}

\begin{theorem}[Bentness of the complete quadratic function]
\lean{FABL.isBent_completeQuadratic}
\label{FABL.isBent_completeQuadratic}
\leanok
\uses{FABL.IsBent, FABL.abs_vectorFourierCoeff_completeQuadratic, FABL.completeQuadratic}
\difficulty{1}
Let $n$ be even. The real sign-valued function obtained by encoding the complete
quadratic form $\sum_{i<j}x_i x_j$ over $\mathbb F_2$ is bent on the $n$-dimensional
binary cube: every Fourier coefficient has absolute value $2^{-n/2}$.
\end{theorem}

\begin{theorem}[Is sign valued complete quadratic]
\lean{FABL.isSignValued_completeQuadratic}
\statementsource{Analysis-of-Boolean-Functions-by-Ryan-ODonnell}{pdf-ffbe9c7b3306-p034-b009}
\label{FABL.isSignValued_completeQuadratic}
\leanok
\uses{FABL.IsSignValued, FABL.completeQuadratic, FABL.completeQuadraticBit, FABL.completeQuadratic_apply}
\difficulty{4}
% dataset-pipeline: stub -- prose pending build_informal_statements.py
Statement pending regeneration.
\end{theorem}

\begin{definition}[Bent direct product]
\lean{FABL.bentDirectProduct}
\statementsource{Analysis-of-Boolean-Functions-by-Ryan-ODonnell}{pdf-ffbe9c7b3306-p152-b001}
\label{FABL.bentDirectProduct}
\leanok
\uses{FABL.F₂Cube}
Given real-valued functions $f$ and $g$ on Boolean cubes of dimensions $n$ and $m$,
respectively, define their direct product on the cube of dimension $n+m$ by splitting
each input into its first $n$ and last $m$ coordinates and multiplying the corresponding
function values.
\end{definition}

\begin{theorem}[Evaluation of the bent direct product]
\lean{FABL.bentDirectProduct_append}
\label{FABL.bentDirectProduct_append}
\leanok
\uses{FABL.F₂Cube, FABL.bentDirectProduct}
\difficulty{1}
For nonnegative integers $n$ and $m$, let $f$ and $g$ be real-valued functions on the
respective binary cubes, and let $x$ and $y$ be points in those cubes. The bent direct
product of $f$ and $g$, evaluated at the concatenation of $x$ and $y$, equals
$f(x)g(y)$.
\end{theorem}

\begin{theorem}[Dot product of concatenated vectors]
\lean{FABL.f₂DotProduct_append}
\leanok
\uses{FABL.F₂Cube, FABL.f₂DotProduct}
\difficulty{3}
For natural numbers $n,m$, vectors $a,x\in\mathbb{F}_2^n$, and vectors
$b,y\in\mathbb{F}_2^m$, the dot product of the concatenations of $a$ with $b$ and $x$
with $y$ equals the sum of the dot products of $a$ with $x$ and $b$ with $y$.
\end{theorem}

\begin{theorem}[Walsh character on concatenated cubes]
\lean{FABL.vectorWalshCharacter_append}
\label{FABL.vectorWalshCharacter_append}
\leanok
\uses{FABL.F₂Cube, FABL.f₂DotProduct_append, FABL.vectorWalshCharacter, FABL.vectorWalshCharacter_apply}
\difficulty{3}
Let $a,x\in\mathbb{F}_2^n$ and $b,y\in\mathbb{F}_2^m$. The vector Walsh character
indexed by the concatenation of $a$ and $b$, evaluated at the concatenation of $x$ and
$y$, equals the product of the vector Walsh characters indexed by $a$ and $b$, evaluated
at $x$ and $y$, respectively.
\end{theorem}

\begin{theorem}[Fourier coefficient of a direct product]
\lean{FABL.vectorFourierCoeff_bentDirectProduct_append}
\label{FABL.vectorFourierCoeff_bentDirectProduct_append}
\leanok
\uses{FABL.F₂Cube, FABL.bentDirectProduct, FABL.bentDirectProduct_append, FABL.vectorFourierCoeff, FABL.vectorFourierCoeff_eq_expect, FABL.vectorWalshCharacter, FABL.vectorWalshCharacter_append}
\difficulty{4}
For all natural numbers $n,m$, real-valued functions $f$ on $\mathbb{F}_2^n$ and $g$ on
$\mathbb{F}_2^m$, and vectors $a\in\mathbb{F}_2^n$, $b\in\mathbb{F}_2^m$, the vector
Fourier coefficient of their direct-product function at the concatenated vector $(a,b)$
equals the product of the coefficient of $f$ at $a$ and the coefficient of $g$ at $b$.
\end{theorem}

\begin{theorem}[Direct products preserve bentness]
\lean{FABL.IsBent.directProduct}
\proofsource{Analysis-of-Boolean-Functions-by-Ryan-ODonnell}{pdf-ffbe9c7b3306-p152-b001}
\label{FABL.IsBent.directProduct}
\leanok
\uses{FABL.F₂Cube, FABL.IsBent, FABL.bentDirectProduct, FABL.vectorFourierCoeff_bentDirectProduct_append}
\difficulty{4}
Let $n,m$ be even, and let $f:\mathbb F_2^n\to\mathbb R$ and $g:\mathbb F_2^m\to\mathbb
R$ be bent. The function on $\mathbb F_2^{n+m}$ that sends a concatenated pair $(x,y)$
to $f(x)g(y)$ is bent.
\end{theorem}

\begin{theorem}[Sign-valued direct products]
\lean{FABL.IsSignValued.directProduct}
\label{FABL.IsSignValued.directProduct}
\leanok
\uses{FABL.F₂Cube, FABL.IsSignValued, FABL.bentDirectProduct}
\difficulty{1}
Let $n,m\ge 0$, and let $f:_2^n	obR$ and $g:_2^m	obR$ satisfy $|f(x)|=|g(y)|=1$ for
every input. Define their direct product on $_2^{n+m}$ by splitting each input into its
first $n$ and last $m$ coordinates and multiplying the corresponding values of $f$ and
$g$. Then this direct product has absolute value $1$ at every input.
\end{theorem}

\begin{theorem}[Fourier coefficient under character multiplication]
\lean{FABL.vectorFourierCoeff_vectorWalshCharacter_mul}
\label{FABL.vectorFourierCoeff_vectorWalshCharacter_mul}
\leanok
\uses{FABL.F₂Cube, FABL.vectorFourierCoeff, FABL.vectorFourierCoeff_eq_expect, FABL.vectorWalshCharacter, FABL.vectorWalshCharacter_mul}
\difficulty{3}
Let $n$ be a natural number, let $f:\mathbb{F}_2^n\to\mathbb{R}$, and let
$a,\gamma\in\mathbb{F}_2^n$. The Fourier coefficient at $\gamma$ of the function
obtained by multiplying $f$ pointwise by the Walsh character indexed by $a$ equals the
Fourier coefficient of $f$ at $a+\gamma$.
\end{theorem}

\begin{definition}[Bent affine modulation]
\lean{FABL.bentAffineModulation}
\statementsource{Analysis-of-Boolean-Functions-by-Ryan-ODonnell}{pdf-ffbe9c7b3306-p152-b003}
\label{FABL.bentAffineModulation}
\leanok
\uses{FABL.F₂Cube, FABL.Sign, FABL.signValue, FABL.vectorWalshCharacter}
Given a sign $\sigma\in\{\pm1\}$, a vector $a$ in the $n$-dimensional Boolean cube,
and a real-valued function $f$ on that cube, define the modulated function at $x$ as the
product of the sign, the Walsh character indexed by $a$ evaluated at $x$, and $f(x)$.
\end{definition}

\begin{theorem}[Vector fourier coeff bent affine modulation]
\lean{FABL.vectorFourierCoeff_bentAffineModulation}
\statementsource{Analysis-of-Boolean-Functions-by-Ryan-ODonnell}{pdf-ffbe9c7b3306-p152-b004}
\label{FABL.vectorFourierCoeff_bentAffineModulation}
\leanok
\uses{FABL.F₂Cube, FABL.Sign, FABL.bentAffineModulation, FABL.signValue, FABL.vectorFourierCoeff, FABL.vectorFourierCoeff_const_mul, FABL.vectorFourierCoeff_vectorWalshCharacter_mul, FABL.vectorWalshCharacter}
\difficulty{3}
% dataset-pipeline: stub -- prose pending build_informal_statements.py
Statement pending regeneration.
\end{theorem}

\begin{theorem}[Affine modulation preserves bentness]
\lean{FABL.IsBent.affineModulation}
\proofsource{Analysis-of-Boolean-Functions-by-Ryan-ODonnell}{
  pdf-ffbe9c7b3306-p152-b003,
  pdf-ffbe9c7b3306-p152-b004
}
\label{FABL.IsBent.affineModulation}
\leanok
\uses{FABL.F₂Cube, FABL.IsBent, FABL.Sign, FABL.bentAffineModulation, FABL.signValue_eq_neg_one_or_one, FABL.vectorFourierCoeff_bentAffineModulation}
\difficulty{4}
Let $n$ be even and let $f:\mathbb F_2^n\to\mathbb R$ be bent. For any sign
$\sigma\in\{\pm1\}$ and any $a\in\mathbb F_2^n$, the function
$x\mapsto\sigma\chi_a(x)f(x)$ is bent, where $\chi_a$ is the Walsh character indexed by
$a$.
\end{theorem}

\begin{theorem}[Affine modulation preserves sign values]
\lean{FABL.IsSignValued.affineModulation}
\label{FABL.IsSignValued.affineModulation}
\leanok
\uses{FABL.F₂Cube, FABL.IsSignValued, FABL.Sign, FABL.abs_vectorWalshCharacter, FABL.bentAffineModulation, FABL.signValue_eq_neg_one_or_one}
\difficulty{4}
Let $n$ be a nonnegative integer, let $f:\\mathbb F_2^n\\to\\mathbb R$ satisfy
$|f(x)|=1$ for every $x$, let $\\sigma\\in\\{-1,1\\}$, and let $a\\in\\mathbb F_2^n$.
Then the function $x\\mapsto\\sigma(-1)^{\\sum_i a_i x_i}f(x)$ also has absolute value
$1$ at every input.
\end{theorem}

\begin{definition}[Affine sign function]
\lean{FABL.affineSignFunction}
\label{FABL.affineSignFunction}
\leanok
\uses{FABL.F₂Cube, FABL.Sign, FABL.bentAffineModulation}
For natural numbers $n,m$, a sign $\sigma\in\mathbb Z^\times$, and $\gamma\in(\mathbb
F_2)^n$, the affine sign function is the function on $(\mathbb F_2)^n$ obtained by
applying the affine sign modulation determined by $\sigma$ and $\gamma$ to the constant
function with value $1$.
\end{definition}

\begin{theorem}[Affine sign function evaluation]
\lean{FABL.affineSignFunction_apply}
\label{FABL.affineSignFunction_apply}
\leanok
\uses{FABL.F₂Cube, FABL.Sign, FABL.affineSignFunction, FABL.bentAffineModulation, FABL.signValue, FABL.vectorWalshCharacter}
\difficulty{1}
For a sign $\sigma\in\{\pm1\}$ and vectors $\gamma,x\in\mathbb F_2^n$, the affine sign
function at $x$ equals the real value of $\sigma$ times the Walsh character indexed by
$\gamma$ evaluated at $x$.
\end{theorem}

\begin{theorem}[Sign encoding of an affine Boolean function]
\lean{FABL.realSignEncodedFunction_affineFunction}
\label{FABL.realSignEncodedFunction_affineFunction}
\leanok
\uses{FABL.F₂Cube, FABL.affineFunction, FABL.affineSignFunction, FABL.affineSignFunction_apply, FABL.binarySign, FABL.f₂DotProduct, FABL.realSignEncodedFunction, FABL.signEncode, FABL.signEncode_add, FABL.signValue, FABL.signValue_signEncode_eq_binarySign, FABL.vectorWalshCharacter, FABL.vectorWalshCharacter_apply, FABL.𝔽₂}
\difficulty{3}
Let $b\in\mathbb F_2$ and $a\in\mathbb F_2^n$. The real sign encoding of the affine
function $x\mapsto b+a\cdot x$ equals the real function $x\mapsto (-1)^b\chi_a(x)$,
where $\chi_a$ is the Walsh character associated with $a$.
\end{theorem}

\begin{theorem}[Affine Sign Functions Are Sign-Valued]
\lean{FABL.isSignValued_affineSignFunction}
\label{FABL.isSignValued_affineSignFunction}
\leanok
\uses{FABL.F₂Cube, FABL.IsSignValued, FABL.Sign, FABL.abs_vectorWalshCharacter, FABL.affineSignFunction, FABL.affineSignFunction_apply, FABL.signValue_eq_neg_one_or_one}
\difficulty{4}
For every natural number $n$, every sign $\\sigma\\in\\{−1,1\\}$, and every vector
$\\gamma\\in\\mathbb F_2^n$, the function on $\\mathbb F_2^n$ given by
$x\\mapsto\\sigma(-1)^{\\sum_{j=1}^n\\gamma_jx_j}$ takes only values of absolute value
$1$.
\end{theorem}

\begin{theorem}[Distance to an affine sign function]
\lean{FABL.relativeHammingDist_affineSignFunction}
\label{FABL.relativeHammingDist_affineSignFunction}
\leanok
\uses{FABL.F₂Cube, FABL.IsSignValued, FABL.Sign, FABL.affineSignFunction, FABL.affineSignFunction_apply, FABL.isSignValued_affineSignFunction, FABL.relativeHammingDist, FABL.signValue, FABL.uniformProbability, FABL.uniformProbability_ne_eq_relativeHammingDist, FABL.vectorFourierCoeff, FABL.vectorFourierCoeff_eq_expect}
\difficulty{5}
Let $f$ be a real-valued function on the binary cube of dimension $n$ such that
$|f(x)|=1$ for every input $x$. For every sign $\sigma\in\{\pm1\}$ and every binary
vector $\gamma$, the relative Hamming distance from $f$ to the affine sign function
determined by $\sigma$ and $\gamma$ is $\frac12-\frac{\sigma}{2}\widehat f(\gamma)$,
where $\widehat f(\gamma)$ is the vector Fourier coefficient of $f$ at $\gamma$.
\end{theorem}

\begin{definition}[Distance to affine sign functions]
\lean{FABL.distanceToAffineSigns}
\label{FABL.distanceToAffineSigns}
\leanok
\uses{FABL.F₂Cube, FABL.Sign, FABL.affineSignFunction, FABL.relativeHammingDist}
For a function $f:\{0,1\}^n\to\mathbb{R}$, define its distance to the affine sign
functions to be the minimum, over all sign vectors and all points of the binary cube, of
the relative Hamming distance from $f$ to the corresponding affine sign function.
\end{definition}

\begin{theorem}[Choosing a sign to realize absolute value]
\lean{FABL.exists_signValue_mul_eq_abs}
\label{FABL.exists_signValue_mul_eq_abs}
\leanok
\uses{FABL.Sign, FABL.signValue, FABL.signValue_neg_one}
\difficulty{3}
For every real number $c$, there is a sign $\sigma\in\{-1,1\}$ such that $\sigma c=|c|$.
\end{theorem}

\begin{theorem}[Distance to affine signs and Fourier norm]
\lean{FABL.distanceToAffineSigns_eq}
\label{FABL.distanceToAffineSigns_eq}
\leanok
\uses{FABL.F₂Cube, FABL.IsSignValued, FABL.Sign, FABL.affineSignFunction, FABL.distanceToAffineSigns, FABL.exists_signValue_mul_eq_abs, FABL.relativeHammingDist, FABL.relativeHammingDist_affineSignFunction, FABL.signValue, FABL.signValue_eq_neg_one_or_one, FABL.spectralInfinityNorm, FABL.vectorFourierCoeff}
\difficulty{5}
Let $n$ be a natural number and let $f$ be a real-valued function on the Boolean cube
with $|f(x)|=1$ at every point. Its normalized Hamming distance to the family of affine
sign functions equals $\frac12-\frac12\|\widehat f\|_\infty$, where $\|\widehat
f\|_\infty$ is the maximum absolute value of its Fourier coefficients.
\end{theorem}

\begin{theorem}[Spectral infinity norm eq]
\lean{FABL.IsBent.spectralInfinityNorm_eq}
\statementsource{Analysis-of-Boolean-Functions-by-Ryan-ODonnell}{pdf-ffbe9c7b3306-p151-b005}
\label{FABL.IsBent.spectralInfinityNorm_eq}
\leanok
\uses{FABL.F₂Cube, FABL.IsBent, FABL.spectralInfinityNorm}
\difficulty{3}
% dataset-pipeline: stub -- prose pending build_informal_statements.py
Statement pending regeneration.
\end{theorem}

\begin{theorem}[Bent bound for affine-sign distance]
\lean{FABL.distanceToAffineSigns_le_bentBound}
\label{FABL.distanceToAffineSigns_le_bentBound}
\leanok
\uses{FABL.F₂Cube, FABL.IsSignValued, FABL.distanceToAffineSigns, FABL.distanceToAffineSigns_eq, FABL.exists_bent_extremal_coefficient, FABL.spectralInfinityNorm, FABL.vectorFourierCoeff}
\difficulty{4}
Let $n$ be an even natural number, and let $f$ be a real-valued function on the Boolean
cube with $|f(x)|=1$ at every point. The normalized Hamming distance from $f$ to the
family of affine sign functions is at most $\frac12-\frac{1}{2\cdot 2^{n/2}}$.
\end{theorem}

\begin{theorem}[Bentness from the spectral maximum]
\lean{FABL.isBent_of_spectralInfinityNorm_eq}
\label{FABL.isBent_of_spectralInfinityNorm_eq}
\leanok
\uses{FABL.F₂Cube, FABL.IsBent, FABL.IsSignValued, FABL.spectralInfinityNorm, FABL.sum_sq_vectorFourierCoeff_eq_one, FABL.vectorFourierCoeff, FABL.𝔽₂}
\difficulty{6}
Let $n$ be an even natural number and let $f$ be a real-valued function on the
$n$-dimensional binary cube such that $|f(x)|=1$ at every input. If the maximum absolute
value of its Fourier coefficients is $2^{-n/2}$, then every Fourier coefficient of $f$
has absolute value $2^{-n/2}$; in particular, $f$ is bent.
\end{theorem}

\begin{theorem}[Bentness and distance to affine signs]
\lean{FABL.isBent_iff_distanceToAffineSigns_eq}
\label{FABL.isBent_iff_distanceToAffineSigns_eq}
\leanok
\uses{FABL.F₂Cube, FABL.IsBent, FABL.IsBent.spectralInfinityNorm_eq, FABL.IsSignValued, FABL.distanceToAffineSigns, FABL.distanceToAffineSigns_eq, FABL.isBent_of_spectralInfinityNorm_eq}
\difficulty{3}
Let $n$ be even, and let $f$ be a sign-valued real function on the $n$-dimensional
binary cube. Then $f$ is bent if and only if its minimum relative Hamming distance to an
affine sign function is $\frac12-\frac12\,2^{-n/2}$.
\end{theorem}

\begin{definition}[Dual frequency under a linear change]
\lean{FABL.bentDualFrequency}
\label{FABL.bentDualFrequency}
\leanok
\uses{FABL.F₂Cube, FABL.𝔽₂}
Let $M$ be an invertible linear map on the length-$n$ vector space over $\mathbb{F}_2$,
and let $\gamma$ be a vector in that space. Define the dual frequency associated to $M$
and $\gamma$ to be the unique vector $\delta$ such that, for every vector $x$,
$\delta\cdot x=\gamma\cdot M^{-1}x$, where the dot product is over $\mathbb{F}_2$.
\end{definition}

\begin{theorem}[Dual frequency under a linear change of variables]
\lean{FABL.f₂DotProduct_bentDualFrequency}
\leanok
\uses{FABL.F₂Cube, FABL.bentDualFrequency, FABL.f₂DotProduct, FABL.𝔽₂}
\difficulty{4}
Let $n$ be a nonnegative integer, let $M$ be an invertible linear map on $br_2^n$, and
let $b gamma,y\in\bbr_2^n$, where $\bbr_2$ is the field with two elements. Define the
dual frequency $M^*\gamma$ by requiring its dot product with $y$ to equal the dot
product of $\gamma$ with $M^{-1}y$. Then the dot product of the bent dual frequency of
$\gamma$ with $y$ equals the dot product of $\gamma$ with $M^{-1}y$.
\end{theorem}

\begin{definition}[Bent linear reindexing]
\lean{FABL.bentLinearReindex}
\statementsource{Analysis-of-Boolean-Functions-by-Ryan-ODonnell}{
  pdf-ffbe9c7b3306-p152-b003,
  pdf-ffbe9c7b3306-p152-b004
}
\label{FABL.bentLinearReindex}
\leanok
\uses{FABL.F₂Cube, FABL.𝔽₂}
Given an invertible linear map $M$ from the $n$-dimensional Boolean cube to itself and a
real-valued function $f$ on that cube, define the reindexed function by $x\mapsto
f(Mx)$.
\end{definition}

\begin{theorem}[Vector fourier coeff bent linear reindex]
\lean{FABL.vectorFourierCoeff_bentLinearReindex}
\statementsource{Analysis-of-Boolean-Functions-by-Ryan-ODonnell}{
  pdf-ffbe9c7b3306-p152-b003,
  pdf-ffbe9c7b3306-p152-b004
}
\label{FABL.vectorFourierCoeff_bentLinearReindex}
\leanok
\uses{FABL.F₂Cube, FABL.bentDualFrequency, FABL.bentLinearReindex, FABL.f₂DotProduct_bentDualFrequency, FABL.vectorFourierCoeff, FABL.vectorFourierCoeff_eq_expect, FABL.vectorWalshCharacter_apply, FABL.𝔽₂}
\difficulty{4}
% dataset-pipeline: stub -- prose pending build_informal_statements.py
Statement pending regeneration.
\end{theorem}

\begin{theorem}[Linear reindexing preserves bentness]
\lean{FABL.IsBent.linearReindex}
\proofsource{Analysis-of-Boolean-Functions-by-Ryan-ODonnell}{
  pdf-ffbe9c7b3306-p152-b003,
  pdf-ffbe9c7b3306-p152-b004
}
\label{FABL.IsBent.linearReindex}
\leanok
\uses{FABL.F₂Cube, FABL.IsBent, FABL.bentLinearReindex, FABL.vectorFourierCoeff_bentLinearReindex, FABL.𝔽₂}
\difficulty{1}
Let $n$ be even, let $f:\mathbb F_2^n\to\mathbb R$ be bent, and let $M:\mathbb
F_2^n\to\mathbb F_2^n$ be an invertible linear map. Then the reindexed function
$x\mapsto f(Mx)$ is bent.
\end{theorem}

\begin{theorem}[Sign values under linear reindexing]
\lean{FABL.IsSignValued.linearReindex}
\label{FABL.IsSignValued.linearReindex}
\leanok
\uses{FABL.F₂Cube, FABL.IsSignValued, FABL.bentLinearReindex, FABL.𝔽₂}
\difficulty{1}
Let $f:_2^n	obR$ satisfy $|f(x)|=1$ for every $x$, and let $M:_2^n	o_2^n$ be an
invertible linear map. Then the function $x	o f(Mx)$ also has absolute value $1$ at
every input.
\end{theorem}

\begin{definition}[Maiorana–McFarland construction]
\lean{FABL.maioranaMcFarland}
\statementsource{Analysis-of-Boolean-Functions-by-Ryan-ODonnell}{pdf-ffbe9c7b3306-p170-b005}
\label{FABL.maioranaMcFarland}
\leanok
\uses{FABL.F₂Cube, FABL.Sign, FABL.f₂CubeBlockEquiv, FABL.innerProductModTwo, FABL.signValue}
Given a function $g$ from the $n$-dimensional Boolean cube to the signs, define a
real-valued function on the Boolean cube of dimension $2n$ by splitting each input into
two $n$-dimensional blocks, taking the inner product modulo two of those blocks, and
multiplying its sign character by the sign value of $g$ evaluated at the second block.
\end{definition}

\begin{theorem}[Maiorana–McFarland formula on joined blocks]
\lean{FABL.maioranaMcFarland_joinF₂CubeBlocks}
\leanok
\uses{FABL.F₂Cube, FABL.Sign, FABL.binarySign, FABL.f₂DotProduct, FABL.innerProductModTwo_joinF₂CubeBlocks, FABL.joinF₂CubeBlocks, FABL.maioranaMcFarland, FABL.signValue}
\difficulty{3}
Let $n$ be a nonnegative integer, let $g$ map the $n$-dimensional cube over $\mathbb
F_2$ to $\{\pm1\}$, and take $x,y\in\mathbb F_2^n$. The Maiorana–McFarland function
evaluated on the joined blocks $(x,y)$ equals
\[(-1)^{\langle x,y\rangle}\,g(y),\]
where the inner product is in $\mathbb F_2$ and is converted to its corresponding sign.
\end{theorem}

\begin{theorem}[Zero sum and equality on the binary cube]
\lean{FABL.f₂Cube_add_eq_zero_iff_eq}
\leanok
\uses{FABL.F₂Cube}
\difficulty{3}
For every natural number $n$ and every two vectors $x,y\in\mathbb{F}_2^n$, their sum is
zero if and only if $x=y$.
\end{theorem}

\begin{theorem}[Orthogonality of binary Walsh characters]
\lean{FABL.expect_innerProductModTwo_mul_vectorWalshCharacter}
\proofsource{Analysis-of-Boolean-Functions-by-Ryan-ODonnell}{
  pdf-ffbe9c7b3306-p152-b006,
  pdf-ffbe9c7b3306-p152-b007
}
\label{FABL.expect_innerProductModTwo_mul_vectorWalshCharacter}
\leanok
\uses{FABL.F₂Cube, FABL.binarySign, FABL.expect_vectorWalshCharacter, FABL.f₂Cube_add_eq_zero_iff_eq, FABL.f₂DotProduct, FABL.vectorWalshCharacter, FABL.vectorWalshCharacter_apply, FABL.vectorWalshCharacter_mul}
\difficulty{4}
Let $n\in\mathbb N$ and let $a,y\in\mathbb F_2^n$. The uniform expectation over
$x\in\mathbb F_2^n$ of $(-1)^{x\cdot y}$ times the Walsh character indexed by $a$ is $1$
if $y=a$ and $0$ otherwise.
\end{theorem}

\begin{theorem}[Fourier coefficient of the Maiorana–McFarland function]
\lean{FABL.vectorFourierCoeff_maioranaMcFarland_joinF₂CubeBlocks}
\leanok
\uses{FABL.F₂Cube, FABL.Sign, FABL.binarySign, FABL.expect_innerProductModTwo_mul_vectorWalshCharacter, FABL.f₂CubeBlockEquiv, FABL.f₂DotProduct, FABL.joinF₂CubeBlocks, FABL.maioranaMcFarland, FABL.maioranaMcFarland_joinF₂CubeBlocks, FABL.signValue, FABL.vectorFourierCoeff, FABL.vectorFourierCoeff_eq_expect, FABL.vectorWalshCharacter, FABL.vectorWalshCharacter_joinF₂CubeBlocks, FABL.𝔽₂}
\difficulty{5}
Let $g$ assign a sign to each point of the $n$-dimensional binary cube, and let $a,b$ be
points of the cube. The vector Fourier coefficient of the Maiorana–McFarland function at
the concatenation of $a$ and $b$ is $2^{-n}$ times the real value of $g(a)$ multiplied
by the Walsh character indexed by $b$ evaluated at $a$.
\end{theorem}

\begin{theorem}[Sign-valued Maiorana–McFarland function]
\lean{FABL.isSignValued_maioranaMcFarland}
\label{FABL.isSignValued_maioranaMcFarland}
\leanok
\uses{FABL.F₂Cube, FABL.IsSignValued, FABL.Sign, FABL.abs_vectorWalshCharacter, FABL.binarySign, FABL.f₂CubeBlockEquiv, FABL.f₂DotProduct, FABL.joinF₂CubeBlocks, FABL.maioranaMcFarland, FABL.maioranaMcFarland_joinF₂CubeBlocks, FABL.signValue_eq_neg_one_or_one, FABL.vectorWalshCharacter_apply}
\difficulty{4}
Let $n$ be a nonnegative integer and let $g$ assign a sign in $\{\pm1\}$ to each vector
in $\mathbb F_2^n$. The function on $\mathbb F_2^{2n}$ that splits an input into blocks
$(x,y)$ and takes the product of $(-1)^{x\cdot y}$ with $g(y)$ has absolute value $1$ at
every input.
\end{theorem}

\begin{theorem}[Maiorana–McFarland bentness]
\lean{FABL.isBent_maioranaMcFarland}
\label{FABL.isBent_maioranaMcFarland}
\leanok
\uses{FABL.F₂Cube, FABL.IsBent, FABL.Sign, FABL.abs_vectorWalshCharacter, FABL.f₂CubeBlockEquiv, FABL.joinF₂CubeBlocks, FABL.maioranaMcFarland, FABL.signValue, FABL.signValue_eq_neg_one_or_one, FABL.vectorFourierCoeff_maioranaMcFarland_joinF₂CubeBlocks}
\difficulty{4}
For every natural number $n$ and every function $g$ from the $n$-dimensional binary cube
to the signs, the function on the $2n$-dimensional binary cube given by the
Maiorana–McFarland construction is bent: every Fourier coefficient has absolute value
$2^{-n}$.
\end{theorem}
